\section{Morris Preorder Traversal}
\label{sec:trees-morris-preorder}

\problemstatement{Return the preorder traversal of a binary tree in $O(1)$
extra space -- no recursion and no explicit stack.}

\begin{patternbox}
Morris preorder is Morris inorder with the ``record'' step moved. Recall
from Section~\ref{sec:trees-rec-traversals} that preorder and inorder are
the same walk with the node recorded at a different moment. In the threaded
walk, that means recording the node \emph{when you create the thread}
(first visit) rather than when you remove it.
\end{patternbox}

\subsection*{Understanding the problem carefully}

The Morris machinery -- thread the predecessor's null right pointer, walk
left, later return and unthread -- visits each node twice when it has a left
child: once on the way in (creating the thread) and once on the way back
(removing it). Inorder records on the second visit; preorder must record on
the \emph{first}. So we output the node's value at the moment we create the
thread (and immediately for nodes with no left child), which places the
node before its subtrees -- exactly preorder.

\begin{ideabox}[Same Threads, Record on the First Visit]
Identical to Morris inorder, with one line moved: when you create a thread
(first time you reach a node with a left child), record the node
\emph{then}, before moving left. Nodes without a left child are still
recorded immediately. When you later remove a thread, do \emph{not} record
again.
\end{ideabox}

\subsection*{Algorithm}

\begin{enumerate}
  \item \textit{curr} $=$ root.
  \item While \textit{curr} is non-null:
  \begin{enumerate}
    \item If \textit{curr} has no left child: record \textit{curr}; move
          right.
    \item Else find predecessor \textit{pred} (rightmost of the left
          subtree, not through a thread):
    \begin{enumerate}
      \item If \textit{pred}.right is null: \textbf{record
            \textit{curr}} (first visit); set \textit{pred}.right $=$
            \textit{curr}; move left.
      \item Else: reset \textit{pred}.right to null; move right (no
            recording).
    \end{enumerate}
  \end{enumerate}
\end{enumerate}

\subsection*{Dry run}

\dryrun{Tree: root $1$; left $2$; right $3$.}

\begin{itemize}
  \item \textit{curr}$=1$, has left. Predecessor $2$.right null $\to$
        record $1$; thread $2\!\to\!1$; move left to $2$.
  \item \textit{curr}$=2$, no left $\to$ record $2$; move right via
        thread to $1$.
  \item \textit{curr}$=1$, predecessor $2$.right $=1$ (thread) $\to$
        remove thread; move right to $3$ (no record).
  \item \textit{curr}$=3$, no left $\to$ record $3$.
  \item Preorder: $1,2,3$.
\end{itemize}

\subsection*{C++ implementation}

\begin{lstlisting}
#include <bits/stdc++.h>
using namespace std;

struct TreeNode {
    int val;
    TreeNode* left;
    TreeNode* right;
    TreeNode(int v) : val(v), left(nullptr), right(nullptr) {}
};

vector<int> morrisPreorder(TreeNode* root) {
    vector<int> out;
    TreeNode* curr = root;
    while (curr != nullptr) {
        if (curr->left == nullptr) {
            out.push_back(curr->val);
            curr = curr->right;
        } else {
            TreeNode* pred = curr->left;
            while (pred->right != nullptr && pred->right != curr)
                pred = pred->right;
            if (pred->right == nullptr) {
                out.push_back(curr->val);    // record on FIRST visit
                pred->right = curr;
                curr = curr->left;
            } else {
                pred->right = nullptr;       // second visit: just unthread
                curr = curr->right;
            }
        }
    }
    return out;
}

int main() {
    TreeNode* root = new TreeNode(1);
    root->left = new TreeNode(2);
    root->right = new TreeNode(3);

    for (int x : morrisPreorder(root)) cout << x << " ";
    cout << "\n";
    return 0;
}
\end{lstlisting}

\begin{complexitybox}
\textbf{Time Complexity.} $O(n)$ -- each edge is walked a constant number
of times. \textbf{Space Complexity.} $O(1)$ -- threads only, tree restored
at the end.
\end{complexitybox}

\begin{takeawaybox}
Once Morris inorder is understood, preorder is a one-line change -- record
on the first visit instead of the second -- reaffirming the chapter's
opening theme that the three depth-first orders differ only in \emph{when}
the node is recorded, threading or not.
\end{takeawaybox}
